\chapter{2006年上海卷（秋文）}
\section{填空题}
\begin{problem}
已知集合 \(A = \{-1, 3, m\}\)，集合 \(B = \{3, 4\}\). 若 \(B \subseteq A\)，则实数 \(m = \)\fillinblank{}.
\end{problem}

\begin{problem}
已知两条直线 \(l_{1}\)：\(ax + 3y - 3 = 0\)，\(l_{2}\)：\(4x + 6y - 1 = 0\). 若 \(l_{1}\parallel l_{2}\)，则 \(a = \)\fillinblank{}.
\end{problem}

\begin{problem}
若函数 \(f(x) = a^x\)（\(a > 0\)，且 \(a \neq 1\)）的反函数的图象过点 \(\left(2,-1\right)\)，则 \(a = \)\fillinblank{}.
\end{problem}

\begin{problem}
计算：\(\lim\limits_{n\to\infty}\frac{n(n^2 + 1)}{6n^3 + 1} =\)\fillinblank{}.
\end{problem}

\begin{problem}
若复数 \(z = (m - 2) + (m + 1)\i\) 为纯虚数（\(\i\)为虚单位），其中 \(m \in \R\)，则 \(|z| = \)\fillinblank{}.
\end{problem}

\begin{problem}
函数 \(y = \sin x \cos x\) 的最小正周期是\fillinblank{}.
\end{problem}

\begin{problem}
已知双曲线的中心在原点，一个顶点的坐标是 \((3,0)\)，且焦距与虚轴长之比为 \(5:4\)，则双曲线的标准方程是\fillinblank{}.
\end{problem}

\begin{problem}
方程 \(\log_{3}(x^{2}-10)=1+\log_{3}x\) 的解是\fillinblank{}.
\end{problem}

\begin{problem}
已知实数 \(x\)，\(y\) 满足 \(\begin{cases} x + y - 3 \geqslant 0,\\ x + 2y - 5 \leqslant 0,\\ x \geqslant 0,\\ y \geqslant 0, \end{cases}\)
则 \(y - 2x\) 的最大值是\fillinblank{}.
\end{problem}

\begin{problem}
在一个小组中有 8 名女同学和 4 名男同学，从中任意地挑选 2 名同学担任交通安全宣传志愿者，那么选到的两名都是女同学的概率是\fillinblank{}（结果用分数表示）.
\end{problem}

\begin{problem}
若曲线 \(|y| = 2^x + 1\) 与直线 \(y = b\) 没有公共点，则 \(b\) 的取值范围是\fillinblank{}.
\end{problem}

\begin{problem}
如图，平面中两条直线 \(l_{1}\) 和 \(l_{2}\) 相交于点 \(O\). 对于平面上任意一点 \(M\)，若 \(p\)，\(q\) 分别是 \(M\) 到直线 \(l_{1}\) 和 \(l_{2}\) 的距离，则称有序非负实数对 \((p, q)\) 是点 \(M\) 的“距离坐标”. 根据上述定义，“距离坐标”是 \((1,2)\) 的点的个数是\fillinblank{}.
\end{problem}

\section{单选题}
\begin{problem}
如图，在平行四边形 \(ABCD\) 中，下列结论中错误的是（\quad）

\bitmapfigure[max width=0.40\linewidth,max totalheight=4.0cm]{img/2006/shanghai_liberal/q13_fig1.png}

\choices
    {\(\overrightarrow{AB} = \overrightarrow{DC}\)}
    {\(\overrightarrow{AD} + \overrightarrow{AB} = \overrightarrow{AC}\)}
    {\(\overrightarrow{AB} - \overrightarrow{AD} = \overrightarrow{BD}\)}
    {\(\overrightarrow{AD} + \overrightarrow{CB} = \overrightarrow{0}\)}
\end{problem}

\begin{problem}
如果 \(a < 0\)，\(b > 0\)，那么，下列不等式中正确的是（\quad）

\choices
    {\(\frac{1}{a} < \frac{1}{b}\)}
    {\(\sqrt{-a} < \sqrt{b}\)}
    {\(a^2 < b^2\)}
    {\(|a| > |b|\)}
\end{problem}

\begin{problem}
若空间中有两条直线，则“这两条直线为异面直线”是“这两条直线没有公共点”的（\quad）

\choices
    {充分非必要条件}
    {必要非充分条件}
    {充分必要条件}
    {既非充分又非必要条件}
\end{problem}

\begin{problem}
如果一条直线与一个平面垂直，那么，称此直线与平面构成一个“正交线面对”. 在一个正方体中，由两个顶点确定的直线与含有四个顶点的平面构成的“正交线面对”的个数是（\quad）

\choices
    {48}
    {18}
    {24}
    {36}
\end{problem}

\section{解答题}
\begin{problem}
已知 \(\alpha\) 是第一象限的角，且 \(\cos \alpha = \frac{5}{13}\)，求 \(\frac{\sin\left(\alpha + \frac{\pi}{4}\right)}{\cos(2\alpha + 4\pi)}\) 的值.
\end{problem}

\begin{problem}
如图，当甲船位于 \(A\) 处时获悉，在其正东方向相距 20 海里的 \(B\) 处有一艘渔船遇险等待营救. 甲船立即前往救援，同时把消息告知在甲船的南偏西 \(30^{\circ}\)，相距 10 海里 \(C\) 处的乙船，试问乙船应朝北偏东多少度的方向沿直线前往 \(B\) 处救援？（角度精确到 \(1^{\circ}\)）

\bitmapfigure[max width=0.44\linewidth,max totalheight=4.6cm]{img/2006/shanghai_liberal/q18_fig1.png}
\end{problem}

\begin{problem}
在直三棱柱 \(ABC - A_{1}B_{1}C_{1}\) 中，\(\angle ABC = 90^{\circ}\)，\(AB = BC = 1\).
\begin{enumerate}
\item 求异面直线 \(B_{1}C_{1}\) 与 \(AC\) 所成角的大小；
\item 若 \(A_{1}C\) 与平面 \(ABC\) 所成角为 \(45^{\circ}\)，求三棱锥 \(A_{1} - ABC\) 的体积.
\end{enumerate}

\bitmapfigure[max width=0.38\linewidth,max totalheight=5.2cm]{img/2006/shanghai_liberal/q19_fig1.png}
\end{problem}

\begin{problem}
设数列 \(\{a_n\}\) 的前 \(n\) 项和为 \(S_n\)，且对任意正整数 \(n\)，\(a_n + S_n = 4096\).
\begin{enumerate}
\item 求数列 \(\{a_n\}\) 的通项公式；
\item 设数列 \(\{\log_2 a_n\}\) 的前 \(n\) 项和为 \(T_n\)，对数列 \(\{T_n\}\)，从第几项起 \(T_n < -509\)？
\end{enumerate}
\end{problem}

\begin{problem}
已知在平面直角坐标系 \(xOy\) 中的一个椭圆. 它的中心在原点，左焦点为 \(F(-\sqrt{3}, 0)\)，且右顶点为 \(D(2, 0)\). 设点 \(A\) 的坐标是 \(\left(1, \frac{1}{2}\right)\).
\begin{enumerate}
\item 求该椭圆的标准方程；
\item 若 \(P\) 是椭圆上的动点，求线段 \(PA\) 中点 \(M\) 的轨迹方程；
\item 过原点 \(O\) 的直线交椭圆于点 \(B\)，\(C\)，求 \(\triangle ABC\) 面积的最大值.
\end{enumerate}
\end{problem}

\begin{problem}
已知函数 \(y = x + \frac{a}{x}\) 有如下性质：如果常数 \(a > 0\)，那么该函数在 \((0, \sqrt{a}]\) 上是减函数，在 \([\sqrt{a}, +\infty)\) 上是增函数.
\begin{enumerate}
\item 如果函数 \(y = x + \frac{2b}{x}\)（\(x > 0\)）在 \((0,4]\) 上是减函数，在 \([4, +\infty)\) 上是增函数，求 \(b\) 的值；
\item 设常数 \(c \in [1,4]\)，求函数 \(f(x) = x + \frac{c}{x}\)（\(1 \leqslant x \leqslant 2\)）的最大值和最小值；
\item 当 \(n\) 是正整数时，研究函数 \(g(x) = x^n + \frac{c}{x^n}\)（\(c > 0\)）的单调性，并说明理由.
\end{enumerate}
\end{problem}
